% ---------------------------------------------------------------------
% Design language for the IJCAR talk.
%
% The chrome is FLoC 26's (floc26.org, sampled from the site assets):
% royal blue #1D3C8A, azulejo gold #F6B655, cream #F4F2F0. The deck's
% one idea --- colour encodes SOURCE vs TARGET --- keeps its own green
% for the source language and takes the brand blue for the target:
%
%   green   = Eunoia   (source)   the pre-rebrand teal
%   blue    = LambdaPi (target)   the brand blue itself
%   gold    = the FLoC accent: chrome (progress bar, rules, tile
%             band) and \alert emphasis (darkened for text); never a
%             language, so it can sit next to either
%   ink     = everything else
%   muted   = supporting detail, comments, grammar delimiters
%
% Chrome comes from the site: solid blue frametitle band with cream
% text (their section bands), gold progress bar (their accents), and
% the azulejo tile strip as a footer on every page --- floc-tiles.png,
% cut seamlessly from the site's own strip asset; the white letter
% tiles hidden in it are theirs. The frame number sits in a cream
% square at the band's right end, echoing those letter tiles.
%
% Consequences, enforced below:
%   - the header band and footer tiles are the only blue chrome; body
%     structure (bullets, alerts) stays ink so blue keeps meaning
%     "LambdaPi" in content
%   - no red anywhere; \alert is bold in darkened FLoC gold
%   - listings are two-tone, not the 5-colour smtlib rainbow
%   - no line numbers
%   - title and divider frames sit on FLoC cream (their hero band);
%     content frames stay near-white
% ---------------------------------------------------------------------

\usepackage{tikz}
\usetikzlibrary{arrows.meta, positioning, fit, backgrounds, shadows.blur}

% ---- measure, margins and vertical rhythm ---------------------------
% At 12pt the line measure was ~56 characters, below the comfortable
% 60-75 band --- which is why the type read as oversized and crowded at
% the same time. 11pt with slightly tighter margins gives ~63.
\setbeamersize{text margin left=8mm, text margin right=8mm}

% Body text indented 24pt from the frametitle's left edge looked like a
% misalignment rather than a hierarchy. Pull the lists back under it.
\setlength{\leftmargini}{12pt}
\setlength{\leftmarginii}{12pt}
\setlength{\leftmarginiii}{12pt}

% One spacing unit, derived from the body size, replacing the ad-hoc
% mix of \vspace{1mm}/{2mm}/{3mm} that was scattered through the deck.
\newlength{\gapunit}
\setlength{\gapunit}{0.6\baselineskip}
\newcommand{\gap}{\vspace{\gapunit}}

% Listings sit in the same rhythm instead of carrying their own skips.
\lstset{aboveskip=\gapunit, belowskip=\gapunit}

% Same for booktabs tables: the top rule carries its own air, so a table
% under a caption line needs no \vspace of its own. (Half a unit, not a
% full one --- the rule is already doing part of the separating.)
\setlength{\abovetopsep}{0.5\gapunit}

% metropolis pins the frametitle's leftskip to its own padding length
% (2.2ex), independent of the text margin, so the title hangs ~12pt to
% the left of every body element. Measured on p.21: title at x=23px,
% body at x=48px. That reads as a misalignment rather than a hierarchy.
% Re-issue the template with the leftskip tied to the text margin,
% leaving the padding length alone --- it also drives the vertical
% struts that give the title block its height.
\makeatletter
\defbeamertemplate{frametitle}{aligned}{%
	\nointerlineskip%
	\begin{beamercolorbox}[%
			wd=\paperwidth,%
			sep=0pt,%
			leftskip=\beamer@leftmargin,%
			rightskip=\beamer@rightmargin,%
		]{frametitle}%
		\metropolis@frametitlestrut@start%
		\insertframetitle%
		\nolinebreak%
		\metropolis@frametitlestrut@end%
	\end{beamercolorbox}%
}
\setbeamertemplate{frametitle}[aligned]
\makeatother

% metropolis implements `progressbar=frametitle` by APPENDING the bar to
% whatever the frametitle template is when \metroset runs; re-issuing the
% template above therefore threw it away, and the deck shipped with no
% progress indicator at all. Re-append it here.
\addtobeamertemplate{frametitle}{}{%
	\usebeamertemplate*{progress bar in head/foot}}

% ---- palette --------------------------------------------------------
% FLoC 26 brand colours, sampled from floc26.org assets. The gold is
% 1.8:1 against white --- chrome, never text.
\definecolor{flocblue}{HTML}{1D3C8A}    % site bands, tile blue
\definecolor{flocamber}{HTML}{F6B655}   % tile gold
\definecolor{floccream}{HTML}{F4F2F0}   % hero background
\colorlet{eocol}{MaterialTeal700}     % Eunoia
\colorlet{eobg}{MaterialTeal50}
\colorlet{lpcol}{flocblue}            % LambdaPi
\colorlet{lpbg}{flocblue!6}
\colorlet{ink}{MaterialGrey900}
\colorlet{muted}{MaterialGrey600}
\colorlet{rule}{MaterialGrey400}

% ---- inline code ----------------------------------------------------
% Prose used bare \texttt, which reaches for the theme's Fira Mono at the
% full body size: measured 10.8pt of ink against the 9pt of the JuliaMono
% inside the boxes, so a passing mention of a tool name shouted louder
% than the generated code the audience is asked to read. One family for
% all code, scaled to the body's x-height so it also behaves in
% frametitles, and the source/target colour rule applies inline too.
%
% \text{} rather than a bare font switch: \kw and friends are used inside
% math displays, where \newfontfamily switches are invalid.
% (\lpcode would clash with LuaTeX's protrusion primitive; the tt suffix
% follows \mtt, which the paper's macros already use for the same job.)
% Ligatures off: JuliaMono is a programming font and renders -> as a
% single arrow glyph, but \kw{->} is the literal Eunoia token. (The
% listings font escapes this because listings sets every character
% separately, which breaks the ligature lookup.)
% -calt as well as the ligature keys: JuliaMono builds -> from contextual
% ALTERNATES, which Ligatures=NoContextual (i.e. -clig) does not touch.
\newfontfamily\lpmonoinline{JuliaMono}[Scale=MatchLowercase,
	Ligatures={NoCommon, NoContextual}, RawFeature={-calt}]
\newcommand{\code}[1]{\text{\lpmonoinline #1}}    % tools, flags, output
\newcommand{\eott}[1]{\text{\lpmonoinline\color{eocol}#1}}   % Eunoia
\newcommand{\lptt}[1]{\text{\lpmonoinline\color{lpcol}#1}}   % LambdaPi

% ---- kill the inherited chrome --------------------------------------
% talk.macros sets structure to indigo and recolours per section; make
% it neutral once, here, so no \setbeamercolor is needed per section.
% Body structure stays ink (not brand blue): in content, blue means
% "LambdaPi", and a bullet must not borrow it.
\setbeamercolor*{structure}{fg=ink}
\setbeamercolor*{normal text}{fg=ink}
\setbeamercolor*{section in head/foot}{fg=muted}
\setbeamercolor*{subsection in head/foot}{fg=muted}

% ---- FLoC chrome ----------------------------------------------------
% Header: the site's solid blue band, cream text. The aligned template
% above already spans \paperwidth, so setting bg is all it takes.
\setbeamercolor*{frametitle}{fg=floccream, bg=flocblue}
% Since the header is orientation (section · subsection), not content,
% it cedes its size: \small against metropolis's \large, and the
% band's padding comes down with it. The font alone barely moves the
% band (10.4mm -> 9.9): the struts are padding-dominated, and the
% padding is a fixed length. Trimming it is safe now --- the aligned
% frametitle template above no longer uses it for leftskip, only the
% struts read it. Measured band: 10.4mm at \large/2.2ex, 8.5mm here.
\setbeamerfont{frametitle}{size=\small}
\makeatletter
\setlength{\metropolis@frametitle@padding}{1.5ex}
\makeatother

% Progress bar under the band: gold fill on a pale blue track, thick
% enough to read as an accent rather than a hairline.
\setbeamercolor{progress bar}{fg=flocamber, bg=flocblue!18}
\makeatletter
\setlength{\metropolis@progressinheadfoot@linewidth}{1.4pt}
\setlength{\metropolis@titleseparator@linewidth}{1.4pt}
\makeatother

% Title page: blue title, gold separator rule (metropolis draws it).
\setbeamercolor{title}{fg=flocblue}
\setbeamercolor{title separator}{fg=flocamber}

% ---- frame headers name the position, not the slide -----------------
% Convention (Ciarán, 2026-07-26): the frametitle is the current
% section, qualified by the current subsection when one is set ---
% orientation lives in the header, specifics live on the slide and in
% the spoken script. Every content frame opens
% \begin{frame}{\secheader}; a \subsection{...} before a run of frames
% qualifies them, \subsection*{} returns to the bare section header.
% The subsection sits in a cream-blue blend: readable on the band,
% clearly subordinate to the section. Emptiness is tested by width ---
% \insertsubsectionhead has no reliably \ifx-able empty form.
\newlength{\subsecheadw}
\newcommand{\secheader}{%
	\insertsectionhead
	\settowidth{\subsecheadw}{\insertsubsectionhead}%
	\ifdim\subsecheadw>0pt
		{\color{floccream!70!flocblue}\ ·\ \insertsubsectionhead}%
	\fi}

% Footer: the FLoC azulejo strip, full-bleed at the bottom of every
% frame (plain ones included --- on the site the strip separates
% sections, so dividers keep it). Drawn on the background canvas, so
% it reserves no layout space; floc-tiles.png is baked at exactly
% \paperwidth : \tilebandh (128:4), so width x height is not a stretch.
\newlength{\tilebandh}
\setlength{\tilebandh}{4mm}
\setbeamertemplate{background canvas}{%
	\usebeamercolor{background canvas}%
	\begin{tikzpicture}
		\useasboundingbox (0,0) rectangle (\the\paperwidth,\the\paperheight);
		\ifbeamercolorempty[bg]{background canvas}{}{%
			\fill[color=bg] (0,0) rectangle (\the\paperwidth,\the\paperheight);}
		\node[anchor=south west, inner sep=0pt] at (0,0)
		{\includegraphics[width=\paperwidth, height=\tilebandh]{floc-tiles}};
	\end{tikzpicture}}

% Frame number: a cream square sitting on the strip's right end, blue
% numeral --- the strip hides white letter tiles, this is the deck's
% own. Replaces the metropolis footline (numbering=fraction), whose
% ~11mm of reserved height the band does not need; the text block only
% grows. Plain frames (title, dividers) get no number, as before.
% 0.3pt taller than the band: at some rasterisations the image and the
% node round to different pixel heights and a hairline of pattern shows
% over the tile; the overhang lands cream-on-white, invisible.
\setbeamertemplate{footline}{%
	\begin{beamercolorbox}[wd=\paperwidth, sep=0pt, right]{footline}%
		\tikz\node[fill=floccream, minimum width=\tilebandh,
			minimum height={\dimexpr\tilebandh+0.3pt\relax}, inner sep=0pt,
			font=\tiny, text=flocblue] {\insertframenumber};%
	\end{beamercolorbox}}

% \alert wears the FLoC gold: with Eunoia back on teal no language
% claims it, so the brand accent is free for emphasis (the site uses
% its orange the same way). The tile gold itself is 1.8:1 on white ---
% an accent, not a text colour --- so alerts take the same hue
% darkened to 5.2:1. Weight comes along: gold alone at \small body
% sizes is not enough of a signal.
\definecolor{alertgold}{HTML}{96610D}
\setbeamercolor*{alerted text}{fg=alertgold}
\setbeamerfont{alerted text}{series=\bfseries}

% \mtt marks Eunoia symbols, so it takes the Eunoia colour (was magenta)
% and the code family (was the theme's mono at body size).
\renewcommand{\mtt}[1]{\eott{#1}}

% "Example." led in MaterialGreen900; make it neutral and quiet.
\renewcommand{\exxample}{{\itshape\color{muted}Example.\ }}

% Grammar annotations were MaterialIndigo500 --- indigo now means
% "LambdaPi", so these must not borrow it.
\renewcommand{\mcomment}[1]{{\color{muted}{\text{#1}}}}

% \smt sets language=smtlib explicitly, which resets to the language's
% own 5-colour keyword styles. Pin it to the two-tone style.
\renewcommand{\smt}[1]{\lstinline[style=eo, basicstyle=\lpmonoinline]{#1}}

% ---- semantic text macros -------------------------------------------
\newcommand{\eoterm}[1]{{\color{eocol}#1}}   % a Eunoia thing, in prose
\newcommand{\lpterm}[1]{{\color{lpcol}#1}}   % a LambdaPi thing, in prose
\newcommand{\eoname}{\eoterm{\textbf{Eunoia}}}
\newcommand{\lpname}{\lpterm{\textbf{LambdaPi}}}
\renewcommand{\LambdaPi}{\lpname}
% τ appears inside math displays, so it takes the inline scale, not the
% listings scale --- at Scale=1.0 it sat a size above the surrounding math.
\renewcommand{\tauop}{\lptt{τ}}

% ---- pass/fail marks (the only pictorial marks in the deck) ---------
% Deliberately uncoloured: teal and indigo mean "Eunoia" and
% "LambdaPi", and a status glyph must not borrow either.
\newcommand{\yes}{{\color{ink}$\checkmark$}}
\newcommand{\no}{{\color{muted}$\times$}}

% A claim item: a list entry marked ✓ rather than bulleted.
% Plain \item[\yes] misaligns. beamer \llap's item labels, so the ink's
% RIGHT edge is pinned \labelsep before the text and a label wider than
% the bullet grows leftwards: the checkmark started 6pt left of the
% frametitle, outside the deck's grid entirely. A negative kern after the
% mark pulls the whole label back into the marker column (measured: ink
% from 57px to 68px, against 74px for the bullets and 63px for the title).
% Takes beamer's overlay spec the way \item does: \claim<+-> ...
\newlength{\claimpull}
\setlength{\claimpull}{4pt}
\makeatletter
\newcommand{\claim@mark}{\yes\hspace*{-\claimpull}}
\def\claim@ov<#1>{\item<#1>[\claim@mark]}
\newcommand{\claim}{\@ifnextchar<{\claim@ov}{\item[\claim@mark]}}
\makeatother

% ---- listings: two-tone, not a rainbow ------------------------------
\lstdefinelanguage{lambdapi}{
	morekeywords={symbol, constant, sequential, injective, rule, with,
			require, open, assert, let, in, notation, unif_rule, builtin,
			opaque, protected, private, inductive, TYPE},
	morekeywords=[2]{Set, Proof, Type},
	alsoletter={:,-,>,|,$},
	sensitive=true,
	morecomment=[l]{//},
	morestring=[b]",
}

% Eunoia source listings.
\lstdefinestyle{eo}{
	language=smtlib,
	basicstyle=\lpmono\small\color{ink},
	keywordstyle=\color{eocol}\bfseries,
	keywordstyle=[2]\color{ink},
	keywordstyle=[3]\color{eocol},
	keywordstyle=[4]\color{ink},
	keywordstyle=[5]\color{eocol}\bfseries,
	commentstyle=\color{muted}\itshape,
	stringstyle=\color{muted},
	numbers=none,   % talk.macros bakes numbers=left into `smtlib` itself
	extendedchars=true,
	showstringspaces=false,
}

% Generated LambdaPi listings.
\lstdefinestyle{lp}{
	language=lambdapi,
	basicstyle=\lpmono\small\color{ink},
	keywordstyle=\color{lpcol}\bfseries,
	keywordstyle=[2]\color{lpcol},
	commentstyle=\color{muted}\itshape,
	stringstyle=\color{muted},
	numbers=none,   % talk.macros bakes numbers=left into `smtlib` itself
	extendedchars=true,
	showstringspaces=false,
}

% No line numbers: nothing in the deck ever refers to one, they were the
% only reason the snippet boxes needed 14pt of left padding against 3pt of
% right, and on a one-line snippet a lone "1" is pure noise.

% SMT-LIB input is neither the source nor the target language. The
% pipeline figure already draws .smt2 as a neutral node; its listing has
% to agree, or teal ends up meaning "Eunoia, and also this other thing".
\lstdefinestyle{smt2}{
	style=eo,
	keywordstyle=\color{ink}\bfseries,
	keywordstyle=[2]\color{ink},
	keywordstyle=[3]\color{ink},
	keywordstyle=[4]\color{ink},
	keywordstyle=[5]\color{ink}\bfseries,
}

\lstset{style=eo}

% ---- snippet boxes --------------------------------------------------
% The paper wraps every code fragment in an eobox/lpbox (preamble.tex).
% Adopted here, with one change: the paper's snippet style is a
% six-colour scheme (pink commands, blue attributes, teal builtins,
% purple types), which would reimport the rainbow this deck removed.
% Instead the BOX carries the language identity via its tint and rule,
% and the code inside keeps the two-tone discipline. The container does
% the semantic work, so the type can stay quiet.
\usepackage{tcolorbox}
\tcbuselibrary{skins}

\colorlet{eoboxbg}{eocol!5}
\colorlet{lpboxbg}{lpcol!5}

% listings inserts trivlist glue inside the box; zero it out.
\newcommand{\snipsetup}{\topsep=0pt \partopsep=0pt \parsep=0pt \parskip=0pt
	\lstset{aboveskip=0pt, belowskip=0pt}}

% Padding is symmetric now that the line-number gutter is gone, and the
% skips are a full \gapunit: at 0.5 they measured 5pt against a 13.6pt
% body line, so prose sat closer to the code frames than to itself --- on
% p16 a single line was 5pt from the box above and 5.4pt from the box
% below, captioning both.
%
% The skips are not free: measured to the text block (not the page
% bottom), the fullest frames had ~1-8pt of slack, so p2, p4 and p20 each
% had to give the space back --- by dropping a \gap that the larger skip
% now provides, and on p2 by joining the continuation lines. Anything
% added to those frames from here needs the same treatment.
\tcbset{snippet/.style={enhanced, arc=0pt, boxrule=0.6pt,
			left=6pt, right=6pt, top=1pt, bottom=1pt,
			before skip=\gapunit, after skip=\gapunit}}

\newenvironment{eobox}{%
	\begin{tcolorbox}[snippet, colback=eoboxbg, colframe=eocol!40]\snipsetup
			}{\end{tcolorbox}}

\newenvironment{lpbox}{%
	\begin{tcolorbox}[snippet, colback=lpboxbg, colframe=lpcol!40]\snipsetup
			}{\end{tcolorbox}}

% Neutral container, for code that belongs to neither language.
\newenvironment{plainbox}{%
	\begin{tcolorbox}[snippet, colback=white, colframe=rule]\snipsetup
			}{\end{tcolorbox}}

% Thin rule between fragments sharing one box (as in the paper).
% \hrule, not \rule-in-a-paragraph: a paragraph line carries
% \baselineskip, which opened almost a full line of air ABOVE the
% hairline and none below --- both split listings read top-heavy.
% Spacing follows the booktabs convention (0.4ex above a midrule,
% 0.65ex below): more GLUE below than above, because the line above
% already contributes its mostly-empty descender band while the caps
% below reach the top of their box --- the INK gaps then read equal.
% Scaled up ~2x from booktabs' table-density values and then tuned on
% a 300dpi render: 0.9/1.05ex measures ~2.5mm of ink gap on each side.
\newcommand{\listsep}[1][eocol]{%
	\par\vspace{0.9ex}%
	{\color{#1!30}\hrule height 0.3pt}%
	\vspace{1.05ex}}

% ---- abstract syntax vocabulary -------------------------------------
% Policy: math displays carry ABSTRACT syntax (metavariables, schemas,
% operators); concrete code only ever appears inside a snippet box.
% Grammar terminals are still concrete tokens, marked with \kw/\attr so
% they read as terminals rather than as pasted code.
\newcommand{\kw}[1]{\text{\lpmonoinline\bfseries\color{eocol}#1}}
\newcommand{\attrib}[1]{\eott{:#1}}
\newcommand{\lpkw}[1]{\text{\lpmonoinline\bfseries\color{lpcol}#1}}
\newcommand{\epar}[1]{%
	\wrapp{muted}{\text{\lpmonoinline(}}{\text{\lpmonoinline)}}{#1}}

% ---- displays -------------------------------------------------------
% $$...$$ centres on the slide and takes TeX's own fixed display skips, so
% each display landed on a left edge of its own --- measured 161pt, 258pt
% and 298pt on p14 alone --- while the frametitle and every body element
% share one edge. Displays join that edge and the \gapunit rhythm.
% Inside a list it flushes to the item's text edge, which is what you
% want there. Also keeps abstract syntax out of \[...\], which cannot be
% left-aligned without fleqn.
\newenvironment{disp}
	{\par\vspace{\gapunit}\noindent$\displaystyle}
	{$\par\vspace{\gapunit}}

% ---- section dividers: typographic, not full-bleed -------------------
% Opening a beamer section has no visual effect (metropolis has no
% headline and sectionpage is off) --- it is kept for the PDF outline.
% FLoC-styled: cream ground (the site's hero band), blue title, gold
% rule; the tile strip arrives via the background canvas. The old
% per-part colour argument is gone --- part colours would fight the
% brand scheme.
% #1 part label, #2 title
\newcommand{\divider}[2]{%
	\section{#2}%
	\begingroup
	\setbeamercolor{background canvas}{bg=floccream}%
	\begin{frame}[plain]
		\vfill
		\hspace{0.12\textwidth}%
		\begin{minipage}{0.8\textwidth}
			{\color{muted}\normalsize #1}\\[1.5mm]
			{\color{flocblue}\Huge\textbf{#2}}\\[3mm]
			{\color{flocamber}\rule{0.3\textwidth}{1.4pt}}
		\end{minipage}
		\vfill
	\end{frame}
	\endgroup}

% ---- the pipeline diagram -------------------------------------------
% Columns carry the language (teal = Eunoia, indigo = LambdaPi), so the
% two eo2lp arrows are parallel and it is visibly the same translation
% applied to a signature and to a proof. Rows carry the frequency:
% the signature is translated once, proofs once per problem.
\newcommand{\pipelinefig}{%
	\resizebox{\textwidth}{!}{%
		\begin{tikzpicture}[
			x=1cm, y=1cm,
			box/.style={draw, rounded corners=1.5pt, align=center,
					font=\footnotesize, inner xsep=4pt, inner ysep=3.5pt,
					minimum height=7.5mm, minimum width=21mm},
			eo/.style={box, draw=eocol, fill=eobg, text=eocol},
			lp/.style={box, draw=lpcol, fill=lpbg, text=lpcol},
			neutral/.style={box, draw=rule, text=ink, minimum width=16mm},
			tool/.style={font=\scriptsize\itshape, text=muted,
					inner sep=1.5pt, fill=white},
			side/.style={font=\scriptsize\itshape, text=muted, anchor=east},
			ar/.style={-{Stealth[length=4.5pt]}, draw=muted, semithick},
		]
		% --- signature row (once) ---
		\node[eo] (cpc)  at (3.2, 1.45) {CPC signature};
		\node[lp] (lcpc) at (7.0, 1.45) {CPC encoding};

		% --- proof row (per problem) ---
		\node[neutral] (smt)  at (0.0, 0) {\texttt{.smt2}};
		\node[eo]      (prf)  at (3.2, 0) {proof script};
		\node[lp]      (lprf) at (7.0, 0) {\texttt{.lp} proof};
		\node[neutral] (ok)   at (10.4, 0) {\yes\ checked};

		\draw[ar] (cpc)  -- node[tool, above] {eo2lp} (lcpc);
		\draw[ar] (smt)  -- node[tool, above] {cvc5} (prf);
		\draw[ar] (prf)  -- node[tool, above] {eo2lp} (lprf);
		\draw[ar] (lprf) -- node[tool, above] {lambdapi} (ok);

		% the encoded signature is what the proof is checked against
		\draw[ar, dashed] (lcpc) -- node[tool, right, xshift=1pt]
		{checked against} (lprf);

			\node[side] at (-0.55, 1.45) {once};
			\node[side] at (-0.55, 0)    {per problem};
		\end{tikzpicture}}}

% ---- the SMT-LIB page stack -----------------------------------------
% Real pages of the standard, fanned about a shared bottom-left corner.
% The artefact makes the "this is a 112-page document" point without
% saying it (it had a "v2.6, 104 pages" caption; cut). Static by
% design: the stack is scene-setting, not an argument being built, so
% it earns no clicks. (The animated variants live in talk-pagestack.tex
% if that ever changes.)
% Requires smtlib.pdf next to the .tex --- the v2.7 reference
% (r2025-07-07); fetch URL in talk-pagestack.tex. PDF page 77 is
% Chapter 5 (Logical Semantics of SMT-LIB Formulas), the part our
% encoding parallels, so it sits in the fan; pdf 55 is a near-blank
% part divider, avoid it.
% 2.3cm: the old 2.0-vs-2.4 ceiling came from the caption line that
% used to sit under the fan; caption gone, the stack earned some of
% the space back (the logic graph beside it gave up width in trade).
\newcommand{\pgstackw}{2.3cm}
\newcommand{\pgstackh}{3.253cm}   % A4 ratio: 2.3 * sqrt(2)
\tikzset{a4page/.style={
	draw=rule, fill=white, line width=0.4pt,
	minimum width=\pgstackw, minimum height=\pgstackh, inner sep=0pt,
	blur shadow={shadow blur steps=10, shadow xshift=1.2pt,
		shadow yshift=-1.6pt, shadow opacity=40}}}
% Two fans, mirror images, centred as a pair with a fixed gutter (they
% originally sat at the ends of the text width bracketing a caption;
% caption cut, gap closed). The left fan pivots on its bottom-left
% corner and opens leftward; the right fan pivots on its bottom-right
% corner and opens rightward, so the two upright top sheets face
% inward. The cover is the left fan's top sheet; every other page is
% just recognisably "dense text". Sheet step is 4 degrees (was 8),
% tightened to pull the deepest sheet's peak corner down --- the peak
% sits at w*sin(step*4) + h*cos(step*4) above the pivot, so most of the
% saving is visual tidiness rather than raw height.
\newcommand{\pagestackfanL}{%
	\begin{tikzpicture}
		\foreach \i/\pg in {5/90, 4/77, 3/45, 2/25} {
			\pgfmathsetmacro{\rot}{(\i-1)*4}
			\node[a4page, rotate around={\rot:(0,0)}, transform shape,
				anchor=south west] at (0,0)
				{\includegraphics[width=\pgstackw, page=\pg]{smtlib.pdf}};
		}
		% Top sheet: the cover ZOOMED, not whole --- drawn entire, its
		% title is ~0.5mm tall. The window is the authors line (the
		% widest band) plus ~20bp of air each side, aspect-matched to
		% the sheet so it fills it exactly, and positioned so the text
		% block (title top to release bottom) is vertically centred in
		% it. Title spans ~75% of the sheet. Trim in bp on the 612x792
		% letter page, measured off a render; re-measure if smtlib.pdf
		% changes revision.
		\node[a4page, anchor=south west] at (0,0)
			{\includegraphics[width=\pgstackw, page=1,
					trim=138bp 173bp 138bp 144bp, clip]{smtlib.pdf}};
	\end{tikzpicture}}
\newcommand{\pagestackfanR}{%
	\begin{tikzpicture}
		\foreach \i/\pg in {5/100, 4/80, 3/60, 2/35, 1/13} {
			\pgfmathsetmacro{\rot}{(\i-1)*4}
			\node[a4page, rotate around={-\rot:(0,0)}, transform shape,
				anchor=south east] at (0,0)
				{\includegraphics[width=\pgstackw, page=\pg]{smtlib.pdf}};
		}
	\end{tikzpicture}}

% ---- the SMT-LIB logic graph ----------------------------------------
% A subset of the standard's logic-inclusion DAG (smt-lib.org draws
% all ~40; the slide needs the shape, not the census). Eight logics:
% the base logics the benchmarks exercise (QF_UF, QF_LIA, QF_AX),
% their joins, and the quantified closures up to AUFLIA. Arrows read
% "is included in". Neutral ink/muted like the pipeline's .smt2 node:
% SMT-LIB is neither the source nor the target language. The lone
% X-crossing in the middle is authentic --- the real graph is a
% thicket, and a planar redraw would misrepresent it politely.
% Sized by the caller: fills \linewidth, so put it in a minipage of
% the width the slide grants it.
\newcommand{\logicgraphfig}{%
	\resizebox{\linewidth}{!}{%
	\begin{tikzpicture}[
		x=1cm, y=1cm,
		logic/.style={font=\footnotesize, text=ink, inner sep=2.5pt},
		ar/.style={-{Stealth[length=4pt]}, draw=muted, thin},
	]
	\node[logic] (qfuf)     at (0, 1.5)     {\code{QF\_UF}};
	\node[logic] (qflia)    at (0, 0)       {\code{QF\_LIA}};
	\node[logic] (qfax)     at (0, -1.5)    {\code{QF\_AX}};
	\node[logic] (qfuflia)  at (2.4, 0.75)  {\code{QF\_UFLIA}};
	\node[logic] (lia)      at (2.4, -0.75) {\code{LIA}};
	\node[logic] (uflia)    at (4.8, 0.75)  {\code{UFLIA}};
	\node[logic] (qfauflia) at (4.8, -0.75) {\code{QF\_AUFLIA}};
	\node[logic] (auflia)   at (7.0, 0)     {\code{AUFLIA}};
	\draw[ar] (qfuf)     -- (qfuflia);
	\draw[ar] (qflia)    -- (qfuflia);
	\draw[ar] (qflia)    -- (lia);
	\draw[ar] (qfax)     to[bend right=10] (qfauflia);
	\draw[ar] (qfuflia)  -- (uflia);
	\draw[ar] (lia)      -- (uflia);
	\draw[ar] (qfuflia)  -- (qfauflia);
	\draw[ar] (uflia)    -- (auflia);
	\draw[ar] (qfauflia) -- (auflia);
	\end{tikzpicture}}}

% ---- proof-tree rule labels (ebproof) --------------------------------
% Quiet, mono, muted: the tree's shape is the point, not the rule names.
\newcommand{\rulelab}[1]{{\color{muted}\scriptsize\code{#1}}}

% ---- the interoperability hub ---------------------------------------
% Ported from talk.tex and restyled: LambdaPi centre in indigo (it is
% the target language), provers in ink, translators as muted mono arrow
% labels, arrows in the pipeline's own style. The PVS arrow was
% unlabelled in the source; left that way rather than guessing a tool.
% 0.88, not 0.92: at 0.92 the hub frame ran 8pt past the text block.
\newcommand{\interopfig}{%
	\resizebox{0.88\textwidth}{!}{%
	\begin{tikzpicture}[
		x=1cm, y=1cm,
		prover/.style={font=\small, text=ink},
		trans/.style={font=\scriptsize, text=muted, fill=white,
				inner sep=1.5pt},
		har/.style={-{Stealth[length=4.5pt]}, draw=muted, semithick},
		hbr/.style={{Stealth[length=4.5pt]}-{Stealth[length=4.5pt]},
				draw=muted, semithick},
	]
	\node[font=\large, text=lpcol] (lp) at (0,0) {\textbf{LambdaPi}};
	\node[prover] (lean) at (-3,2)  {Lean};
	\node[prover] (coq)  at (0,2)   {Rocq};
	\node[prover] (hol)  at (3,2)   {HOL};
	\node[prover] (isa)  at (5,0)   {Isabelle};
	\node[prover] (agda) at (3,-2)  {Agda};
	\node[prover] (pvs)  at (0,-2)  {PVS};
	\node[prover] (k)    at (-3,-2) {$\mathbb{K}$ framework};
	\node[prover] (mat)  at (-5,0)  {Matita};
	\draw[har] (lean) -- (lp) node[trans, midway, sloped, above]
		{\code{lean2dk}};
	\draw[hbr] (coq)  -- (lp) node[trans, midway, sloped, above]
		{\code{vodk}};
	\draw[hbr] (hol)  -- (lp) node[trans, midway, sloped, above]
		{\code{hol2dk}};
	\draw[har] (isa)  -- (lp) node[trans, midway, sloped, above]
		{\code{isabelle\_dk}};
	\draw[har] (agda) -- (lp) node[trans, midway, sloped, above]
		{\code{agda2dk}};
	\draw[har] (pvs)  -- (lp);
	\draw[har] (k)    -- (lp) node[trans, midway, sloped, above]
		{\code{KaMeLo}};
	\draw[hbr] (mat)  -- (lp) node[trans, midway, sloped, above]
		{\code{Krajono}};
	\end{tikzpicture}}}
